% Thorne & Campolattaro 1967, ApJ 149, 591
% Page 15 (p. 605) transcription

The entire theory remains rather academic until one uses it to compute numerically
the complex normal modes for particular, physically interesting stellar models. Such
computations should not be difficult with modern techniques and machines; and they
should yield results of considerable physical interest (cf.\ \S~1).

In subsequent papers in this series, we hope to move in both of the above directions
and also to analyze the dipole ($l = 1$) pulsations which were omitted from consideration
here.

The analysis reported in this paper was guided in part by the Havener (1966) and
Havener-Thorne (1965) analysis of the radial pulsations of cylindrical configurations of
perfect fluid, where gravitational radiation also comes into play. We thank Mr.\ Havener
for helpful discussions. We also gratefully acknowledge several valuable discussions with
Professors John A.\ Wheeler and S.\ Chandrasekhar, without which our analysis would
have been much less complete; and we thank Dr.\ L.\ Edelstein and Dr.\ C.\ V.\ Vishvesh\-wara
for communicating to us in advance of publication some of the results of their
work on perturbations of the Schwarzschild geometry. Finally, we thank the National
Science Foundation (K.\ S.\ T.) and NATO (A.\ C.) for postdoctoral fellowships, and
Princeton University for hospitality, during the spring of 1966 when this investigation
was being initiated.

\appendix

\section*{APPENDIX A \\ THE SPLIT INTO TENSORIAL SPHERICAL HARMONICS \\ AND THE REGGE-WHEELER CHOICE OF GAUGE}

\textit{(cf.\ Regge and Wheeler 1957)}

At the beginning of the analysis we expand the arbitrary metric of an equilibrium configuration in spherical harmonics. All quantities which transform as scalar fields under rotations are
expanded in scalar spherical harmonics, $Y^l_m(\theta, \phi)$, which have \emph{even parity}, $\pi = (-1)^l$.
All fields which transform as vectors with respect to rotations are expanded in vector spherical
harmonics, which are of two types: \emph{even-parity} harmonics ($\pi = (-1)^l$)
%
\begin{equation}
\psi^l_{Mj} = \partial_j Y^l_M(\theta, \phi)
\label{eq:A1a}
\tag{A1a}
\end{equation}
%
and \emph{odd-parity} harmonics ($\pi = (-1)^{l+1}$)
%
\begin{equation}
\phi^l_{Mj} = \epsilon_j{}^{kl}\, \partial_k Y^l_M(\theta, \phi) \,.
\label{eq:A1b}
\tag{A1b}
\end{equation}
%
Here $j$ and $k$ run over $x^3 = \theta$ and $x^4 = \phi$, and $\epsilon_j^{\ kl}$ is the antisymmetric tensor with respect to
rotations
%
\begin{equation}
\epsilon_{34} = -1/\sin\theta \,, \qquad
\epsilon_{43} = \sin\theta \,, \qquad
\epsilon_{33} = \epsilon_{44} = 0 \,.
\label{eq:A2}
\tag{A2}
\end{equation}
%
All fields which transform as tensors with respect to rotations are expanded in tensor spherical
harmonics, which have \emph{even parity}
%
\begin{equation}
\Psi^l_{Mjk} = Y^l_M \gamma_{jk} \,, \qquad
\hat{\Psi}^l_{Mjk} = \partial_j \partial_k Y^l_M \,,
\label{eq:A3a}
\tag{A3a}
\end{equation}
%
or \emph{odd parity}
%
\begin{equation}
\chi^l_{Mjk} = \tfrac{1}{2}\bigl(\epsilon_j{}^{kl}\, \partial_k \psi^l_{Ml}
+ \epsilon_k{}^{lm}\, \partial_l \psi^l_{Mj}\bigr) \,.
\label{eq:A3b}
\tag{A3b}
\end{equation}
%
Here $\gamma_{jk}$ is the metric for rotations
%
\begin{equation}
\gamma_{33} = 1 \,, \qquad
\gamma_{23} = \gamma_{32} = 0 \,, \qquad
\gamma_{44} = \sin^2\theta \,;
\label{eq:A4}
\tag{A4}
\end{equation}
%
and the slash in equation (A3a) denotes covariant differentiation with respect to the metric $\boldsymbol{\gamma}$.
